\chapter{Drugie twierdzenie Gödla i twierdzenie Löba}
\label{ch:second-incompleteness}

\section{Czy teoria może dowieść własnej niesprzeczności?}
Pierwsze twierdzenie Gödla, udowodnione w rozdziale~\ref{ch:first-incompleteness},
wyznacza ograniczenia rozstrzygania zdań w efektywnie aksjomatyzowalnych
teoriach arytmetycznych. Teraz pytamy o jedno szczególne zdanie:
czy teoria może dowieść, że nie istnieje dowód sprzeczności
\emph{w tej samej teorii}? Aby pytanie miało ścisły sens,
ustalamy sposób kodowania dowodów.

\begin{definicja}[Formalna niesprzeczność]\label{def:con-t}
Niech $T$ będzie efektywnie aksjomatyzowalną teorią klasycznej logiki
pierwszego rzędu w języku arytmetyki, zawierającą $\mathrm{PA}$.
Niech $\operatorname{Prf}_T(p,x)$ będzie standardowym predykatem
z Lematu~\ref{res:7:7}, a
$\operatorname{Bew}_T(x):=\exists p\,\operatorname{Prf}_T(p,x)$.
Ustalmy zdanie $\bot:=(0=S(0))$ i jego kod $b:=\ulcorner\bot\urcorner$.
\emph{Zdaniem formalnej niesprzeczności} teorii $T$ nazywamy
\[
\operatorname{Con}_T:=
\neg\operatorname{Bew}_T(\overline b)
\qquad\text{(równoważnie: }\neg\Box\bot\text{)}.
\]
Jest to zdanie języka arytmetyki, a nie wypowiedź
metamatematyczna. Ponieważ $\operatorname{Bew}_T$ jest $\Sigma_1$,
zdanie $\operatorname{Con}_T$ jest $\Pi_1$.
\end{definicja}

\begin{przyklad}[Co kwantyfikuje zdanie niesprzeczności?]
Po rozwinięciu definicji $\operatorname{Con}_T$ otrzymujemy
$\forall p\,\neg\operatorname{Prf}_T(p,\overline b)$.
Dla pojedynczej standardowej liczby $n$ zdanie
$\neg\operatorname{Prf}_T(\overline n,\overline b)$ wyklucza jedynie
jeden kod. Nawet gdy dla każdego standardowego $n$ mamy osobny
dowód tej negacji, nie oznacza to, że istnieje jeden dowód
zdania z kwantyfikatorem $\forall p$. To rozróżnienie będzie
kluczowe przy modelach niestandardowych.
\end{przyklad}

\begin{lemat}[Znaczenie niesprzeczności w modelu standardowym]
\label{lem:con-standard}
Dla $T$ z Definicji~\ref{def:con-t} zachodzi
$\mathfrak N\models\operatorname{Con}_T$ wtedy i tylko wtedy,
gdy $T$ jest niesprzeczna.
\end{lemat}
\begin{proof}
Z definicji wybranego predykatu $\operatorname{Prf}_T$
wynika, że $\mathfrak N\models
\operatorname{Bew}_T(\overline b)$ dokładnie wtedy,
gdy istnieje rzeczywista, skończona dedukcja zdania
$\bot$ z aksjomatów $T$. Ponieważ $T$ zawiera
$\mathrm{PA}$, dowodzi $\neg\bot$.
Jeżeli więc dowodzi także $\bot$, jest sprzeczna.
Odwrotnie, jeżeli $T$ jest sprzeczna, to w klasycznej
logice pierwszego rzędu dowodzi dowolnego zdania,
w szczególności $\bot$ (zasada eksplozji).
Brak dowodu $\bot$ jest zatem równoważny
metamatematycznej niesprzeczności $T$.
\end{proof}

\begin{uwaga}[Niesprzeczność nie jest poprawnością semantyczną]
Prawdziwość $\operatorname{Con}_T$ w $\mathfrak N$
nie oznacza, że wszystkie aksjomaty $T$ są prawdziwe
w $\mathfrak N$. Niesprzeczna teoria może mieć modele,
ale nie mieć modelu standardowego. W dalszym ciągu
nie zakładamy poprawności semantycznej $T$ ani
$\omega$-niesprzeczności.
\end{uwaga}

\section{Od dowodliwości zdania do dowodliwości sprzeczności}
W Twierdzeniu~\ref{thm:xi-hbl-verified} wykazano,
że ustalony predykat dowodliwości spełnia warunki
(D1)--(D3) z Definicji~\ref{def:hilbert-bernays-loeb}.
Wewnętrzne operacje na kodach uzasadniają te warunki;
nie wolno dodawać nieudowodnionej implikacji
$\Box\varphi\to\varphi$.

\begin{lemat}[Wewnętrzne rozumowanie o dowodach]
\label{lem:second-box-g}
Niech $T$ będzie efektywnie aksjomatyzowalną teorią
rozszerzającą $\mathrm{PA}$, a $\Box$ jej standardowym
predykatem dowodliwości. Jeśli zdanie arytmetyczne
$G$ spełnia $T\vdash G\leftrightarrow\neg\Box G$,
to $T\vdash\Box G\to\Box\bot$.
\end{lemat}
\begin{proof}
Z kierunku $G\to\neg\Box G$ równoważności diagonalnej
mamy $T\vdash G\to\neg\Box G$.
Warunek (D1) stosujemy do \emph{całej gotowej implikacji},
a potem (D2):
\begin{equation}\label{eq:second-imp-neg}
T\vdash\Box(G\to\neg\Box G),
\qquad T\vdash\Box G\to\Box\neg\Box G.
\end{equation}
Niezależnie (D3) daje
\begin{equation}\label{eq:second-pos}
T\vdash\Box G\to\Box\Box G.
\end{equation}
Z tautologii $\neg\Box G\to(\Box G\to\bot)$
na mocy (D1) otrzymujemy
$T\vdash\Box(\neg\Box G\to(\Box G\to\bot))$.
Zastosowanie (D2) dla zdań $\neg\Box G$ oraz
$\Box G\to\bot$ daje
$T\vdash\Box\neg\Box G\to\Box(\Box G\to\bot)$.
Drugie zastosowanie (D2), tym razem dla
$\Box G$ oraz $\bot$, daje
$T\vdash\Box(\Box G\to\bot)\to
(\Box\Box G\to\Box\bot)$.
Składając te dwie implikacje, uzyskujemy
\[
T\vdash\Box\neg\Box G\to
(\Box\Box G\to\Box\bot).
\]
Wewnątrz dowodu w $T$ załóżmy tymczasowo
$\Box G$. Z~\eqref{eq:second-imp-neg}
wynika $\Box\neg\Box G$,
z~\eqref{eq:second-pos} wynika $\Box\Box G$,
a z ostatniej implikacji --- $\Box\bot$.
Po usunięciu założenia $\Box G$ za pomocą
twierdzenia o dedukcji otrzymujemy
$T\vdash\Box G\to\Box\bot$, zgodnie z tezą.
Nigdzie nie korzystaliśmy z $\Box G\to G$.
\end{proof}

\begin{twierdzenie}[Drugie twierdzenie Gödla o niezupełności]
\label{thm:godel-second}
Niech $T$ będzie niesprzeczną, efektywnie
aksjomatyzowalną teorią klasycznej logiki pierwszego
rzędu w języku arytmetyki, zawierającą $\mathrm{PA}$.
Dla standardowego predykatu dowodliwości
z Definicji~\ref{def:con-t} zachodzi
$T\nvdash\operatorname{Con}_T$.
\end{twierdzenie}
\begin{proof}
Z Lematu diagonalnego~\ref{res:8:1}, zastosowanego
do formuły $\neg\operatorname{Bew}_T(x)$,
otrzymujemy zdanie $G$, dla którego
$Q\vdash G\leftrightarrow\neg\Box G$.
Ponieważ $T\supseteq\mathrm{PA}\supseteq Q$,
ta sama równoważność jest dowodliwa w $T$.
Na mocy Lematu~\ref{lem:second-box-g} mamy
$T\vdash\Box G\to\Box\bot$.
Kontrapozycja, definicja $\operatorname{Con}_T$
oraz drugi kierunek równoważności diagonalnej dają
\begin{equation}\label{eq:con-implies-g}
T\vdash\operatorname{Con}_T\to\neg\Box G,
\qquad T\vdash\neg\Box G\to G,
\qquad T\vdash\operatorname{Con}_T\to G.
\end{equation}
Przypuśćmy, że $T\vdash\operatorname{Con}_T$.
Wówczas~\eqref{eq:con-implies-g} daje $T\vdash G$.
Warunek (D1) zastosowany do tego \emph{rzeczywistego}
dowodu $G$ daje $T\vdash\Box G$.
Jednocześnie kierunek $G\to\neg\Box G$ punktu
stałego daje $T\vdash\neg\Box G$, wbrew
niesprzeczności $T$. Zatem
$T\nvdash\operatorname{Con}_T$.
Nie potrzebowaliśmy $\omega$-niesprzeczności.
\end{proof}

\begin{przyklad}[Arytmetyka Peana]
Jeśli $\mathrm{PA}$ jest niesprzeczna, to
$\mathrm{PA}\nvdash\operatorname{Con}_{\mathrm{PA}}$.
Nie oznacza to, że formalne zdanie o jej
niesprzeczności jest fałszywe: z
Lematu~\ref{lem:con-standard} wynika wówczas
$\mathfrak N\models\operatorname{Con}_{\mathrm{PA}}$.
Wniosek dotyczy ograniczeń dowodzenia
\emph{wewnątrz tej samej teorii}.
\end{przyklad}

\section{Twierdzenie Löba}
Drugie twierdzenie ma następujące uogólnienie:
jeśli teoria potrafi wykazać, że dowodliwość
pewnego zdania pociąga samo to zdanie, to już
potrafi udowodnić to zdanie bez tego założenia.

\begin{twierdzenie}[Twierdzenie Löba]\label{thm:loeb}
Niech $T\supseteq Q$ będzie teorią klasycznej
logiki pierwszego rzędu w języku arytmetyki.
Niech $\operatorname{Bew}_T(x)$ będzie ustaloną
formułą arytmetyczną, a $\Box\varphi$
skrótowym zapisem
$\operatorname{Bew}_T(\overline{\ulcorner\varphi\urcorner})$
dla każdego zdania $\varphi$. Załóżmy, że dla wszystkich
zdań $\varphi,\psi$ spełnione są warunki (D1)--(D3)
z Definicji~\ref{def:hilbert-bernays-loeb}.
Wtedy dla każdego zdania $\varphi$ zachodzi
\[
T\vdash\Box\varphi\to\varphi
\quad\Longrightarrow\quad T\vdash\varphi.
\]
Nie zakładamy niesprzeczności $T$.
\end{twierdzenie}
\begin{proof}
Ustalmy zdanie $\varphi$ i załóżmy
$T\vdash\Box\varphi\to\varphi$.
Z Lematu diagonalnego~\ref{res:8:1}
dla formuły
$\psi(x):=(\operatorname{Bew}_T(x)\to\varphi)$
otrzymujemy zdanie $L$ takie, że
\begin{equation}\label{eq:loeb-fixed}
T\vdash L\leftrightarrow(\Box L\to\varphi).
\end{equation}
Zmienna $x$ jest liczbowa; zapis $\Box x$
nie byłby formułą naszego języka.

\emph{Krok 1: wykazujemy $T\vdash\Box L\to\Box\varphi$.}
Z pierwszego kierunku~\eqref{eq:loeb-fixed}
mamy $T\vdash L\to(\Box L\to\varphi)$.
Korzystając z (D1) i (D2), otrzymujemy
\[
T\vdash\Box(L\to(\Box L\to\varphi)),
\qquad
T\vdash\Box L\to\Box(\Box L\to\varphi).
\]
Warunek (D2) zastosowany do zdań $\Box L$
oraz $\varphi$ daje
$T\vdash\Box(\Box L\to\varphi)\to
(\Box\Box L\to\Box\varphi)$.
Z kolei (D3) daje $T\vdash\Box L\to\Box\Box L$.
Załóżmy pomocniczo $\Box L$.
Pierwsza implikacja dostarcza
$\Box(\Box L\to\varphi)$, a (D3) --- $\Box\Box L$.
Podstawiając je do powyższej instancji (D2),
otrzymujemy $\Box\varphi$. Usunięcie
założenia pomocniczego daje
\begin{equation}\label{eq:loeb-first}
T\vdash\Box L\to\Box\varphi.
\end{equation}

\emph{Krok 2: dowodzimy $L$.}
Z założenia twierdzenia
$T\vdash\Box\varphi\to\varphi$
i~\eqref{eq:loeb-first}
wynika $T\vdash\Box L\to\varphi$.
Drugi kierunek~\eqref{eq:loeb-fixed}
daje $T\vdash(\Box L\to\varphi)\to L$.
Stosując modus ponens, otrzymujemy $T\vdash L$.

\emph{Krok 3: kończymy dowód.}
Dopiero teraz zastosowanie (D1)
do gotowego dowodu $L$ daje $T\vdash\Box L$.
Na mocy~\eqref{eq:loeb-first} uzyskujemy
$T\vdash\Box\varphi$, a z hipotezy
$T\vdash\Box\varphi\to\varphi$
otrzymujemy $T\vdash\varphi$.
\end{proof}

\begin{wniosek}[Drugie twierdzenie z twierdzenia Löba]
\label{cor:godel-from-loeb}
Niech $T$ będzie niesprzeczną teorią klasyczną
zawierającą $Q$, a jej predykat $\Box$
spełnia (D1)--(D3). Wtedy
$T\nvdash\neg\Box\bot$.
\end{wniosek}
\begin{proof}
Przypuśćmy, że $T\vdash\neg\Box\bot$.
Tautologia klasyczna
$\neg\Box\bot\to(\Box\bot\to\bot)$
daje $T\vdash\Box\bot\to\bot$.
Z Twierdzenia~\ref{thm:loeb}, zastosowanego
do $\varphi=\bot$, wynika $T\vdash\bot$,
wbrew niesprzeczności.
\end{proof}

\section{Modele, które „widzą” dowód sprzeczności}
\begin{twierdzenie}[Niestandardowe modele teorii z zaprzeczeniem niesprzeczności]
\label{thm:models-not-con}
Niech $T\supseteq\mathrm{PA}$ będzie niesprzeczną,
efektywnie aksjomatyzowalną teorią klasycznej logiki
pierwszego rzędu. Wtedy
$T\cup\{\neg\operatorname{Con}_T\}$ jest
niesprzeczna i ma model. Żaden model tego rozszerzenia
nie jest izomorficzny ze standardową strukturą $\mathfrak N$.
\end{twierdzenie}
\begin{proof}
Gdyby $T\cup\{\neg\operatorname{Con}_T\}$
była sprzeczna, twierdzenie o dedukcji dawałoby
$T\vdash\neg\operatorname{Con}_T\to\bot$,
a klasyczna logika --- $T\vdash\operatorname{Con}_T$.
Przeczy to Twierdzeniu~\ref{thm:godel-second}.
Rozszerzenie jest więc niesprzeczne.
Z twierdzenia Gödla o pełności,
Twierdzenie~\ref{res:5:3}, wynika istnienie modelu.
Z Lematu~\ref{lem:con-standard} mamy
$\mathfrak N\models\operatorname{Con}_T$,
więc $\mathfrak N$ nie spełnia
$\neg\operatorname{Con}_T$.
Ponieważ prawdziwość zdań jest zachowywana
przez izomorfizmy struktur, żaden model
rozszerzenia nie jest izomorficzny z $\mathfrak N$.
\end{proof}

\begin{uwaga}[Skąd bierze się niestandardowy świadek?]
Niech $\mathfrak M\models
T\cup\{\neg\operatorname{Con}_T\}$.
Z definicji zaprzeczenia niesprzeczności
istnieje $p\in M$ takie, że
$\mathfrak M\models\operatorname{Prf}_T(p,\overline b)$.
Nie może to być żadna standardowa liczba
$\overline n^{\mathfrak M}$.
Ponieważ $T$ jest niesprzeczna,
dla każdego standardowego $n$ relacja
$\operatorname{Prf}_T(n,b)$ jest fałszywa.
Z rozstrzygania zamkniętych instancji
$\Delta_0$ w $Q$
(Lemat~\ref{lem:incomp-proof-check}) wynika
$Q\vdash\neg\operatorname{Prf}_T(\overline n,\overline b)$.
Model $\mathfrak M$ spełnia $Q$, więc żaden
$\overline n^{\mathfrak M}$ nie jest świadkiem.
Element $p$ jest niestandardowym kodem w
rozumieniu $\mathfrak M$, a nie rzeczywistą,
skończoną dedukcją sprzeczności w metateorii.
\end{uwaga}

\begin{uwaga}[Zależność od predykatu dowodu]
Wniosek $T\nvdash\operatorname{Con}_T$
dotyczy określonej formalizacji dowodliwości,
dla której udowodniono (D1)--(D3).
Formuły zgodne w standardowym modelu nie muszą
być dowodliwie równoważne w $T$.
Nie należy bez dodatkowego dowodu stosować
tego twierdzenia do arbitralnego predykatu,
w szczególności do predykatu Rossera
z Definicji~\ref{def:rosser-predicate}.
\end{uwaga}

\section{Granice interpretacji}
\begin{uwaga}[Czego twierdzenia nie mówią?]
Nie udowodniliśmy niesprzeczności $\mathrm{PA}$
ani tego, że zdanie
$\operatorname{Con}_{\mathrm{PA}}$
nie może zostać udowodnione w teorii silniejszej.
Jeśli teoria $U$ dowodzi takiego zdania, jej
argument jest względny wobec aksjomatów $U$.
Drugie twierdzenie Gödla nie podważa poprawnych
algorytmów weryfikacji dowodów ani nie uzasadnia
twierdzeń o nieomylności ludzkiego rozumowania.
\end{uwaga}
